\subsection{Troncature par irréductibles et contrôle du reste}
\label{ix-add:corte-primos}

La profondeur d'occupation $K$ et la borne de l'irréductible $Y$ définissent
des approximations différentes de la même représentation arithmétique. Les bornes
précédentes conservent tous les irréductibles et tronquent la profondeur. On
maintient maintenant les contributions locales complètes et on restreint les modes
produits par le sélecteur à $p\leq Y$, avec $Y\geq2$ entier. On définit
\[
 B_1(Y)=\gamma_0+\sum_{p\leq Y}
       \left(\log(1-p^{-1})+p^{-1}\right),\qquad
 C_2(Y)=\prod_{2<p\leq Y}\left(1-\frac1{(p-1)^2}\right).
\]
Le produit vide vaut un. La partie finie $\gamma_0$ est le même objet que celui désigné par $\gamma$ dans la formule de Meissel--Mertens ; son évaluation avec reste est développée dans la section~\ref{lectura:manuscrito-04}.

\HMTresultado{Proposition}{Encadrements par coupure des irréductibles}{ix-add:prime-cut-bound}
Pour tout entier $Y\geq2$,
\[
 B_1(Y)-\frac1{2Y}\leq B_1\leq B_1(Y),\qquad
 \frac{Y-1}{Y}C_2(Y)\leq C_2\leq C_2(Y).
\]
\textbf{Démonstration.}
Pour chaque $p\geq2$, la contribution soustraite de $B_1(Y)$ est
\[
 d_p=-\log(1-p^{-1})-p^{-1}
     =\sum_{k\geq2}\frac1{kp^k},\qquad
 0<d_p\leq\frac1{2p(p-1)}.
\]
La dernière inégalité utilise $1/k\leq1/2$ et la somme géométrique. Par comparaison positive et télescopage,
\[
 0\leq B_1(Y)-B_1=\sum_{p>Y}d_p
 \leq\frac12\sum_{n>Y}\frac1{n(n-1)}=\frac1{2Y}.
\]
Chaque facteur de la queue de $C_2$ appartient à $(0,1)$. Inclure tous les
entiers $n>Y$ diminue le produit. Pour $M>Y$,
\[
 \prod_{n=Y+1}^{M}\left(1-\frac1{(n-1)^2}\right)
 =\prod_{n=Y+1}^{M}\frac{n(n-2)}{(n-1)^2}
 =\frac{M(Y-1)}{Y(M-1)}.
\]
La limite est $(Y-1)/Y$. La convergence du produit sur les nombres premiers et la comparaison précédente donnent le second encadrement. \hfill$\square$

Lorsque $Y$ augmente, $B_1(Y)$ et $C_2(Y)$ décroissent vers leurs valeurs respectives.
Les largeurs garanties sont $1/(2Y)$ et $C_2(Y)/Y$. On peut prendre l'intersection des encadrements par
profondeur et par irréductible pour raffiner l'évaluation ;
chacun conserve son paramètre et sa démonstration. Le produit $C_2$
correspond au secteur $p>2$ ; le facteur local du mode deux maintient la
normalisation $2C_2$ de la série singulière complète.
